\chapter{A sharp boundary for hidden classical explanations}
\label{rr:reset-boundary}

A response reversal can establish that a declared experimental model is
inadequate without identifying its replacement. Hidden detector states,
imperfect preparation, and uncharacterized crossing rates can all matter.
This chapter turns two of those alternatives into quantitative requirements.
Its main result interpolates sharply between a protected regenerative interface
and unrestricted coarse observation. It is a theorem about specified stochastic
models, with no assumption of historical priority or empirical confirmation.

\section{The measured quantity and the physical assumptions}

Let $R$ be a finite continuous-time Markov generator and $B$ a detector region.
An initial distribution $\nu_A$ is supported outside $B$ and $\nu_B$ inside $B$.
For any initial distribution $\nu$, define
\[
 U_\nu(T)=\int_0^T\Pr_\nu(X_s\in B)\,ds,
 \qquad H(T)=U_{\nu_B}(T)-U_{\nu_A}(T).
\]
The normalized contrast $H(T)/T$ is directly observable by preparing either
initial distribution and reading the same binary detector at an independently
chosen uniform time in $[0,T]$. This formulation requires neither stationary
preparation nor a current measurement. Preparation errors and detector errors
must still be propagated into the contrast interval.

In the stationary selective-channel experiment developed earlier, the same
quantity enters through
\[
 D_B(T)=[\pi e^{TR}](B)-\pi(B)=\gamma H(T),
 \qquad \pi R=\gamma(\nu_B-\nu_A).
\]
Thus the present result can also interpret a finite-pulse response, provided the
forcing identity and any division by $\gamma$ have been separately justified.
The direct preparation experiment and this response-based experiment are
separate access contracts.

We assume uniform \emph{crossing-hazard} bounds
\begin{equation}
 \sum_{y\in B}R_{xy}\leq a\quad(x\notin B),\qquad
 \sum_{y\notin B}R_{xy}\leq b\quad(x\in B).
 \label{rr:hazards}
\end{equation}
These are bounds on every relevant hidden state, not averages inferred from
steady occupancy. Rates of motion within either detector region need not be
bounded. In particular, a proposed classical hidden-state model is not entitled
to use only the dissipative jump rates of a Lindblad description as its hazard
caps: that substitution would itself need physical justification.

Let $\tau$ be the first entrance into $B$ starting from $\nu_A$. Write $\mu$ for
the possibly defective distribution of $\tau$ and $\zeta_u$ for the entrance
state distribution conditional on $\tau=u$. A reset-defect bound $\eta$ means
\begin{equation}
 \operatorname{TV}(\zeta_u,\nu_B)\leq\eta,
 \qquad \mu\text{-almost every }u\leq T,
 \qquad 0\leq\eta\leq1.
 \label{rr:reset-defect}
\end{equation}
The conditioning on entrance time is substantive. An unconditional average
entrance distribution does not establish this condition. When $\nu_B$ is the
single state $j$, the defect at time $u$ is precisely the conditional probability
of entering through a different hidden state. A nonzero defect can arise in a
perfectly Markovian classical model; it is not synonymous with non-Markovianity.

\section{A sharp hazard--reset tradeoff}

For $c\geq0$ set
\[
 A_c(T)=\begin{cases}(1-e^{-cT})/c,&c>0,\\T,&c=0,\end{cases}
\]
and define
\begin{align}
 G_{a,b}(T)&=\int_0^T e^{-bs-a(T-s)}\,ds,\label{rr:G}\\
 K_{a,b}(T)&=T-A_a(T)-A_b(T)+G_{a,b}(T).\label{rr:K}
\end{align}
Here $G_{a,a}(T)=Te^{-aT}$ and, for $a\ne b$,
$G_{a,b}(T)=(e^{-aT}-e^{-bT})/(b-a)$.

\begin{theorem}[Sharp coarse-detector response bound]
\label{rr:sharp-reset}
Under \eqref{rr:hazards} and \eqref{rr:reset-defect},
\begin{equation}
 H(T)\ \geq\ G_{a,b}(T)-\eta K_{a,b}(T)
 =(1-\eta)G_{a,b}(T)+\eta\{A_a(T)+A_b(T)-T\}.
 \label{rr:sharp-bound}
\end{equation}
The lower bound is attained by a finite classical Markov model for every
$a,b>0$, $T>0$, and $\eta\in[0,1]$.
\end{theorem}

\begin{proof}
Write $p_B(s)=\Pr_{\nu_B}(X_s\in B)$ and
$S(u)=\Pr_{\nu_A}(\tau>u)$. The strong Markov property gives
\[
 U_{\nu_A}(T)=\int_{[0,T]}U_{\zeta_u}(T-u)\,\mu(du).
\]
Adding and subtracting the same expression with $\zeta_u$ replaced by $\nu_B$
yields the exact decomposition
\begin{align*}
 H(T)&=G_{\rm actual}(T)-\mathcal R(T),\\
 G_{\rm actual}(T)&=\int_0^T p_B(s)S(T-s)\,ds,\\
 \mathcal R(T)&=\int_{[0,T]}
       [U_{\zeta_u}(T-u)-U_{\nu_B}(T-u)]\,\mu(du).
\end{align*}
The hazard caps imply $p_B(s)\geq e^{-bs}$ and $S(u)\geq e^{-au}$ by
considering trajectories that have not yet left their initial region.
Consequently $G_{\rm actual}(T)\geq G_{a,b}(T)$.

For every initial state $x\in B$,
$A_b(v)\leq U_x(v)\leq v$. Total variation therefore bounds the reset term by
\[
 \mathcal R(T)\leq
 \eta\,\mathbb E[(T-\tau)_+-A_b((T-\tau)_+)].
\]
The function $v-A_b(v)$ is nonnegative and increasing. Since the entrance
hazard is at most $a$, $\tau$ stochastically dominates an exponential variable
of rate $a$; mass at $\tau=\infty$ contributes zero. For $a>0$, the last
expectation is at most
\begin{align*}
 \int_0^T ae^{-au}[T-u-A_b(T-u)]\,du
 &=\int_0^T[1-e^{-a(T-s)}][1-e^{-bs}]\,ds\\
 &=K_{a,b}(T).
\end{align*}
The formula extends continuously to $a=0$ or $b=0$, when $K=0$.
These estimates prove the inequality.

For attainment, use four states $i,j,k,\ell$, detector $B=\{j,k\}$,
$\nu_A=\delta_i$, $\nu_B=\delta_j$, and the only transitions
\[
 i\longrightarrow j:\ a(1-\eta),\qquad
 i\longrightarrow k:\ a\eta,\qquad
 j\longrightarrow\ell:\ b.
\]
States $k$ and $\ell$ are absorbing. The entrance time is exponential with
rate $a$ and its independent entrance law is
$(1-\eta)\delta_j+\eta\delta_k$. Moreover,
$U_j(v)=A_b(v)$ and $U_k(v)=v$. Every inequality above is an equality.
\end{proof}

The small-time defect penalty is
$K_{a,b}(T)=abT^3/6+O(T^4)$, whereas
$G_{a,b}(T)=T-(a+b)T^2/2+O(T^3)$.
The bounds are therefore sensitive to the time available for both entrance and
exit. The sharpness claim concerns the stated class, not the precision attainable
at fixed energy on an unspecified apparatus.

For $0<\eta<1$, the same bound is the infimum within irreducible models
compatible with the stationary selective-channel experiment. To see this,
add $j\to i$, $k\to j$, and $\ell\to i$ at rate $\varepsilon$, and replace
$j\to\ell$ by $b-\varepsilon$. For sufficiently small positive $\varepsilon$,
the chain is irreducible, both hazard caps and the entrance-reset law remain
unchanged, and its kernel tends to the attaining kernel. Double the two
$i\leftrightarrow j$ rates to define the native generator $Q$. It has a positive
stationary law $\pi$, and suppressing that channel by one half returns $R$.
The response identity is then exact. The native current tends to zero in this
limit; approximate sharpness does not confer a finite-sample advantage for
measuring a shrinking current.

\section{Comparison preparations that already occupy the detector}

The outside-support condition cannot be inferred from a preparation's name or
from an analogy with a qubit basis. Suppose instead that the comparison state is
\[
 \rho=p\nu_{\rm in}+(1-p)\nu_{\rm out},\quad
 \operatorname{supp}\nu_{\rm in}\subseteq B,\quad
 \operatorname{supp}\nu_{\rm out}\subseteq B^c,
\]
with $\operatorname{TV}(\nu_{\rm in},\nu_B)\leq\eta_0$. Let the first-entrance
reset defect from $\nu_{\rm out}$ be bounded by $\eta_1$.
\begin{proposition}[Sharp mixed-preparation extension]
Under the same hazard caps,
\[
 U_{\nu_B}(T)-U_\rho(T)\geq
 (1-p)\{G_{a,b}(T)-\eta_1 K_{a,b}(T)\}
 -p\eta_0\{T-A_b(T)\}.
\]
\end{proposition}
\begin{proof}
Decompose the contrast into its inside and outside contributions. Apply the
previous theorem to the outside contribution and
$U_{\nu_B}-U_{\nu_{\rm in}}\geq-\eta_0(T-A_b)$ to the inside contribution.
The same four-state extremizer, with entrance mixture parameter $\eta_1$ and
$\nu_{\rm in}=(1-\eta_0)\delta_j+\eta_0\delta_k$, attains the result.
\end{proof}
With no reset assumptions, $\eta_0=\eta_1=1$, the sharp bound simplifies to
\[
 U_{\nu_B}(T)-U_\rho(T)\geq A_b(T)+(1-p)A_a(T)-T.
\]
This expression also follows directly from conditional crossing-hazard survival
and therefore has the broader classical-process scope described next. If only
an upper bound on $p$ is known, use that upper bound. For $p=1$ there need not be
any positive protection interval. A preparation with one-half detector
occupancy must not be analyzed using the $p=0$ formula.

\section{An explanation bound without a protected reset assumption}

Setting $\eta=1$ removes the reset assumption entirely:
\begin{equation}
 H(T)\geq A_a(T)+A_b(T)-T.
 \label{rr:unrestricted-coarse}
\end{equation}
This special case has a still broader interpretation. Consider any classical
binary process with the stated initial detector value and with uniformly bounded
conditional crossing intensities: at most $a$ from outside to inside and at most
$b$ in the reverse direction, under every admissible history. Survival before
the first crossing obeys the same exponential bounds. At each time $s$,
\[
 \Pr_{\nu_B}(B,s)\geq e^{-bs},\qquad
 \Pr_{\nu_A}(B,s)\leq1-e^{-as}.
\]
Integration proves \eqref{rr:unrestricted-coarse}. No finite hidden dimension,
Markov property, stationarity, or specified first-entrance reset is needed for
this case. The hypothesis of finite uniformly bounded crossing intensities
remains essential; deterministic jumps and unbounded boundary chatter do not
satisfy it.

With a common cap $a,b\leq\lambda$, define
\[
 F(z)=\frac{2(1-e^{-z})}{z}-1,\qquad F(0)=1.
\]
Then $H(T)/T\geq F(\lambda T)$. The function $F$ decreases strictly from $1$
to $-1$. Its zero is $z_0=1.593624260\ldots$, defined exactly by
$z_0=2(1-e^{-z_0})$. Hence an integrated contrast cannot become negative before
$T=z_0/\lambda$, regardless of the number of hidden classical states. The
supplied rational certificates place $z_0$ strictly between $1.593$ and $1.594$.

This result can be used without pretending that a classical hazard cap has
already been experimentally certified. If $c_U\in(-1,1)$ is a calibrated upper
confidence bound on $H(T)/T$, every classical explanation with finite conditional
crossing intensities must have
\begin{equation}
 \lambda\geq\frac{F^{-1}(c_U)}{T}.
 \label{rr:activity-inverse}
\end{equation}
This is a lower bound on the peak crossing intensity required of that explanation.
It becomes a rejection test only when an independently justified upper cap is
smaller. Observing a negative contrast by itself neither excludes all classical
models nor establishes a thermodynamic energy cost for their activity.

\section{What finite data can establish}

Under calibrated caps $a,b$ with $K_{a,b}(T)>0$, an upper confidence bound $H_U$
on $H(T)$ implies
\begin{equation}
 \eta\geq\max\left\{0,
       \frac{G_{a,b}(T)-H_U}{K_{a,b}(T)}\right\}.
 \label{rr:defect-inverse}
\end{equation}
A value exceeding one is inconsistent with the declared bounded-hazard
classical contract. A positive value below one states a minimum required reset
defect. It does not identify which microscopic state or physical mechanism
caused that defect. If $K=0$, the response bound is independent of $\eta$ and
one must test that bound directly instead of dividing by zero.

For a simple direct protocol, collect $N$ independent binary observations from
each initial preparation, choosing the observation time independently and
uniformly in $[0,T]$ for every record. Let $\widehat c$ be the difference of the
two empirical means in the order $B$ minus $A$. Weighted Hoeffding gives
\[
 \Pr\{H(T)/T>\widehat c+u\}\leq e^{-Nu^2}
\]
under the ideal model. If separately calibrated per-arm record laws differ in
total variation by at most $\epsilon_B$ and $\epsilon_A$, use
$c_U=\widehat c+u+\epsilon_B+\epsilon_A$. A joint calibration failure allowance
is added once. The bound assumes independent records and predetermined sample
counts; adaptive stopping requires a valid sequential procedure. The exact
binomial procedure in the statistical chapter offers a complementary encoded
single-bit implementation.

The software implements the sharp bound and its inversions with rational
exponential enclosures. An output certificate establishes a conditional
mathematical statement, not the correctness of the supplied hardware
calibration. The hazard-inversion helper gives both endpoints of a certified
root bracket. It does not report the numerical root as an exact physical rate.

\section{A sharp four-state requirement without a hazard calibration}
\label{rr:four-states}

The geometry of a binary detector gives an additional explanation bound that
does not require a kinetic cap or a protected reset interface.
\begin{theorem}[Minimum classical state count for a negative integrated contrast]
\label{rr:dimension}
Consider a common finite, time-homogeneous classical Markov generator, a fixed
binary detector region $B$, and initial distributions supported respectively
inside and outside $B$. If the hidden state space has at most three states,
then $H(T)>0$ for every finite $T>0$. Thus $H(T)<0$ requires at least four
classical Markov states. Four states suffice to realize any prescribed value
$H(T)/T\in(-1,0)$ at a prescribed positive $T$.
\end{theorem}
\begin{proof}
Both detector regions are nonempty. With at most three states, one of them is a
singleton. If $B=\{j\}$, the first-hitting return argument gives
\[
 U_j(T)-U_{\nu_A}(T)
 =\int_0^T P_j(X_s=j)\Pr_{\nu_A}(\tau_j>T-s)\,ds>0.
\]
The inequality is strict because a finite generator permits positive no-jump
probabilities at every finite time. If instead $B^c=\{i\}$, complement the
detector:
\[
 U^B_{\nu_B}(T)-U^B_i(T)
 =U^{B^c}_i(T)-U^{B^c}_{\nu_B}(T)>0,
\]
and apply the same argument. This covers arbitrary mixtures within the allowed
initial supports.

For sufficiency, use the four-state attaining model with $\eta=1$ and
$a=b=\lambda$. Its contrast is $F(\lambda T)$, which decreases continuously
from $1$ to $-1$. Every specified contrast in $(-1,0)$ therefore determines a
unique positive $\lambda$ and is attained exactly.
\end{proof}

The matching claim concerns one integrated contrast. Four states need not
reproduce the full time-resolved record, multiple intervention responses, or
a complete quantum process. State count also does not determine stationary
Shannon entropy, physical memory cost, or thermodynamic dissipation.

A qubit supplies a direct mathematical comparison. Under unitary evolution
with Hamiltonian $gX$, initial states $|1\rangle$ and $|0\rangle$, and detector
$|1\rangle\langle1|$, the contrast is
\[
 \frac{H(T)}{T}=\frac{\sin(2gT)}{2gT}.
\]
At $T=3\pi/(4g)$ it equals $-2/(3\pi)$. A two-dimensional quantum system can
therefore produce a contrast that requires at least four states in the stated
classical Markov representation. This is a conditional state-count separation,
not a device-independent proof of quantumness or a claim that classical
four-state models are inadequate for that scalar observation.

\begin{proposition}[A conservative preparation-robust state-count witness]
\label{rr:dimension-robust}
Allow a fixed stochastic binary detector with response $f(x)\in[0,1]$ and
arbitrary initial preparations $\nu_B,\nu_A$. Write
\[
 m_B=\nu_B f,\quad m_A=\nu_A f,\quad
 C_f(T)=\frac1T\int_0^T(\nu_B-\nu_A)e^{sR}f\,ds.
\]
Write $C_0=m_B-m_A$. Every homogeneous classical model with at most three
states satisfies the directly observable bound
\begin{equation}
 \boxed{\left|C_0-C_f(T)\right|\leq1.}
 \label{rr:dimension-drop}
\end{equation}
In particular,
\begin{equation}
 C_f(T)\geq m_B-m_A-1.
 \label{rr:dimension-noise}
\end{equation}
\end{proposition}
\begin{proof}
First consider a deterministic detector region $D$. Conditional preparations
inside and outside $D$ have nonnegative integrated contrast by the preceding
theorem. Replacing the actual preparations by these conditional preparations
changes a normalized occupation expectation by at most
$\nu_B(D^c)+\nu_A(D)$ in total. If a conditional preparation is undefined, this
sum is at least one and the bound follows from the range $[-1,1]$.
Thus the normalized contrast for $D$ is at least
$-\nu_B(D^c)-\nu_A(D)$.

Use the layer-cake representation
$f(x)=\int_0^1\mathbf1_{\{f(x)>u\}}\,du$ and integrate that bound over the
threshold detector regions. The two integrated impurity terms are
$1-\nu_B f$ and $\nu_A f$, proving \eqref{rr:dimension-noise}. Swapping the
two preparations proves the opposite bound and hence
\eqref{rr:dimension-drop}.
\end{proof}

The bound in \eqref{rr:dimension-drop} is sharp as a supremum for the
three-state class. Let a reference state be absorbing with detector response
$c\in[0,1]$. Let the comparison state have detector response zero and jump at
rate $\lambda$ to a third absorbing state with detector response one. Then
$C_0=c$, $C_f(T)=c-1+A_\lambda(T)/T$, and
$C_0-C_f(T)=1-A_\lambda(T)/T\to1$ as $\lambda\to\infty$.
The more specific impurity form is not claimed optimal for every separately
specified pair of initial means.

With simultaneous confidence
bounds $C_U$, $m_{B,L}$, and $m_{A,U}$, the condition
\[
 C_U<m_{B,L}-m_{A,U}-1
\]
excludes every classical model with at most three homogeneous Markov states
under the fixed-detector contract. It requires no statement that either
preparation has exactly zero hidden impurity. The same observable bound applies to every independent nonincreasing
observation-clock density starting at zero. The clock chapter represents such
a density as a mixture of uniform clocks, so
\[
 C_w=\int C_f(T)\,\mu(dT),\qquad
 |C_0-C_w|\leq\int|C_0-C_f(T)|\,\mu(dT)\leq1.
\]
A clock chosen using the evolving hidden trajectory is not covered by this
argument.

The directly observable score $C_0-C_f(T)$ needs no assumption of exactly
pure hidden preparations. It requires the same detector response function at
baseline and at the random observation time, and the same preparations in the
corresponding arms. Common fixed readout noise is already covered by
$f\in[0,1]$. If a detector is permitted to reverse its labels between baseline
and readout, a motionless two-state system can give $C_0=1$ and $C_f(T)=-1$;
readout stability is therefore essential.

A common known delay $d$ before both baseline and random-time observations does
not invalidate \eqref{rr:dimension-drop}. Replace the preparations by
$\nu_B e^{dR}$ and $\nu_A e^{dR}$, which are allowed to be arbitrary. The same
argument covers an independent common delay distribution used in every arm,
by averaging these effective preparations; the delay must also be independent
of the observation-clock draw. Different delay distributions,
changed generators, or different detector calibrations require separate error
bounds.

The four-state equal-hazard construction can match both averaged arm
probabilities of a unital qubit with orthogonal preparations and a projective
binary detector: they are $A_\lambda(T)/T$ and $1-A_\lambda(T)/T$, equivalently
$(1+C)/2$ and $(1-C)/2$ once $F(\lambda T)=C$. This statement applies when the
random timestamps are discarded. Keeping the timestamp, intervening within
records, or comparing several clocks yields additional data that the scalar
matching argument does not reproduce.

Time homogeneity is essential. A three-state classical counterexample starts
in $j$ or $i$ with detector $\{j\}$. First allow $j\to k$ at rate ten for time
$1/10$; next allow $i\to j$ at rate ten for time $1/10$; then freeze all rates
for time one. The integrated contrast over the whole interval is
\[
 H(6/5)=\frac{19}{10e}-\frac9{10}<0.
\]
Thus an external control clock, unmodeled memory, or preparation-dependent
kinetics can invalidate the state-count conclusion. The theorem compares
common autonomous Markov representations; these restrictions must be checked
before interpreting the result as a hidden-state requirement.

\section{Protected graph interfaces beyond one entrance gateway}

For a known directed support graph, additional protections are available without
numerical rate identification. Write $U_x(T)$ for the occupation integral from
state $x$. Two elementary first-hitting comparisons hold for every rate choice:
\begin{enumerate}
\item If $x\in B$ and every directed path from $x$ to $B^c$ visits $y$, then
      $U_x(T)\geq U_y(T)$.
\item If $y\notin B$ and every directed path from $y$ to $B$ visits $x$, then
      $U_x(T)\geq U_y(T)$.
\end{enumerate}
For the first, all time before hitting $y$ is spent in $B$, and
$U_y(T)-U_y(T-s)\leq s$. For the second, no time before hitting $x$ is spent in
$B$, and $U_x$ is increasing. The strong Markov property gives each comparison,
including trajectories that never reach the gateway.

These inequalities compose transitively. A reachability test after deleting a
candidate gateway checks every primitive, and transitive closure supplies a
verifiable chain of comparisons. For example, on the graph
$i\to k$, $k\to\ell$, $j\to k$, $\ell\to k$ with $B=\{j,k\}$,
neither $j$ as the entrance gateway from $i$ nor $i$ as the exit gateway from
$j$ is valid. Nevertheless $U_j\geq U_k\geq U_i$. A closed detector region is
also protected by the exit condition. The implemented preorder is a sufficient
certificate. A complete graph characterization of universal occupation order
is not claimed, and failure to find a certificate is not evidence of a
counterexample.

\section{Why exact hidden structure cannot be certified by finite records}

A finite calibration experiment cannot uniformly rule out arbitrarily weak
unmodeled transitions. Suppose extra transitions have total intensity at most
$\delta$ at every state and under every allowed control. Coupling the original
and perturbed processes until the first extra jump gives, for records with total
pathwise horizon at most $L$,
\[
 \operatorname{TV}(P_{\rm original},P_{\rm perturbed})
 \leq1-e^{-\delta L}\leq\delta L.
\]
This statement includes adapted controls that obey the same intensity bound;
any detector or downstream processing preserves the inequality. Consequently,
a finite-record rule that accepts the original structural explanation with
probability $p$ accepts a sufficiently small perturbation with probability at
least $p-\delta L$.

The issue can affect all-time protection, not merely exact graph labels. Take
$i\to k$ at rate one and $j\to\ell$ at rate $\delta$, with
$B=\{j,k\}$ and $k,\ell$ absorbing. At $\delta=0$,
$H(T)=1-e^{-T}>0$ for every positive $T$. At every $\delta>0$,
\[
 H(T)=A_\delta(T)+A_1(T)-T
\]
eventually becomes negative, while finite-horizon records approach those of
the protected model as $\delta\downarrow0$. Structural physical knowledge,
a quantitative tolerance, or a positive separation assumption is therefore
needed. Taking additional finite records alone cannot establish exact absence
of all arbitrarily weak hidden alternatives.

The practical output is an explicit choice: certify a physical interface,
report a quantitative minimum reset defect, or report the hidden peak crossing
intensity any competing classical explanation would require. Each result says
what the evidence establishes and what remains open.
